\chapter*{全书符号与约定}
\addcontentsline{toc}{chapter}{全书符号与约定}

实数域与复数域分别记为 $\RR$ 与 $\CC$，微分统一写作正体 $\dd$。谱理论部分的复 Hilbert 空间内积约定对第一变量线性，例如 $\langle u,v\rangle_w=\int u\overline v\,w\dd x$；涉及其他内积时在首次出现处明确其类型。微分几何中恒等映射记为 $\IdMap$；挠率张量及其二形式分量统一记为 $T$ 与 $T^a$，曲率二形式记为 $\Omega^a{}_b$。这样把 $\Theta$ 保留给统计部分的参数空间，避免几何与统计在相邻模块中复用同一大写希腊字母。有边界 Hodge 理论中外法向量统一记为 $\nu$，absolute/relative Laplacian 分别记为 $\Delta_A,\Delta_R$；主丛联络一形式与曲率分别记为 $\omega$ 与 $\Omega$，其作用空间不同于 moving-frame 指标化曲率 $\Omega^a{}_b$，首次出现时均显式注明类型。

数理统计中单次观测空间统一写为 $(\mathcal X,\mathcal A)$，参数空间写为 $\Theta$，模型族写为 $\mathcal P=\{P_\theta:\theta\in\Theta\}$；$\rho$ 不用作统计参数。$\GammaD(\alpha,\lambda)$ 始终采用 shape--rate 参数化，其密度含因子 $\lambda^\alpha\ee^{-\lambda x}$。收敛记号、经验测度与 likelihood 约定在统计部分首次出现时固定，后续章节沿用同一含义。
